\section{Trabalhos Relacionados}

\color{black}

Trabalhos recentes demonstram a superioridade de abordagens baseadas em IA sobre análise clássica de
sinais em ambientes ruidosos. Na vertente supervisionada, \cite{Le2025} propuseram modelos avaliados nas
bases MIMII e DCASE usando características de alta frequência (Espectrogramas de Mel e MFCCs) combinadas
a XGBoost e GRU, com resultados expressivos em ambientes industriais complexos. Corroborando a
aplicabilidade no chão de fábrica, \cite{Uzel2026} aplicaram CNNs leves sobre espectrogramas de Mel para
diagnóstico de motores elétricos, demonstrando latências na casa dos milissegundos em dispositivos de
borda.

No regime não supervisionado, o desafio central reside no \textit{Domain Shift} e na escassez de amostras
anômalas --- o paradigma \textit{First-Shot} do DCASE~2025~\cite{Nishida2025}. \cite{Bi2026} integram
aprimoramento espectral e \textit{Frequency-Gated Attention} para lidar com a ausência de amostras de
falha. Soluções competitivas do DCASE~2025 priorizam otimização contrastiva latente e modelagem
adversarial para robustecer fronteiras de decisão frente a variações operacionais~\cite{Jiang2025,
Emon2025}. O alinhamento ao protocolo \textit{First-Shot} do DCASE~2025 foi estabelecido por
\cite{Harada2023} como referência formal de generalização cross-domain.

Para contornar o peso computacional do Deep Learning, modelos leves e DSP ganham notoriedade.
\cite{Zamorano2025} atestaram a Transformada Wavelet de Pacote (WPT) para identificação de transientes
mecânicos ultracurtos. \cite{Ali2025} utilizaram PCA acoplada a Autoencoders Adversariais para acelerar
a detecção em séries temporais multivariadas. \cite{Zhang2026} introduziram o Tiny-AST, modelo compacto
($\sim$1M parâmetros) de alta performance em condições adversas de ruído. Num estudo de referência para
Edge Computing, \cite{Khamaisi2025} validaram em hidrelétricas reais que modelos estatísticos rasos
(K-Means, OC-SVM, Distância de Mahalanobis) atingem o melhor balanço entre tempo de treinamento e
precisão (AUC~$>$~0,96) quando alicerçados por engenharia de características refinada.

\textbf{Lacuna identificada:} A literatura carece de um estudo que: (a)~documente formalmente e
quantifique \textit{o que foi descartado} e \textit{por que}; (b)~proponha uma arquitetura dual
supervisionada/não-supervisionada com decisão limiar estatisticamente fundamentada; e
(c)~valide com múltiplas fontes e métricas DCASE-alinhadas. Este trabalho preenche essa lacuna.

\color{black}
